\chapter{Main Chapter}
\label{ch:mainchapter}

Main chapter. This can be duplicated several times.

% comment out the following lines if you do not need a table of content for each chapter
\newpage
\minitoc % table of content for each chapter
\newpage

%%%%%%%%%%%
% Figures  
%%%%%%%%%%%
\section{Figures}

Example figure (see Figure~\ref{fig:Bristollogo}).

\subsection{Single figure}
\begin{figure}[h]
    \centering
    \includegraphics[width=0.8\columnwidth]{images/Logo - colour PNG.png}
    \caption{Bristol Logo.}
    \label{fig:Bristollogo}
\end{figure}

\subsection{Figures with sub captions}

\begin{figure}[h]
    \centering
    \begin{subfigure}[h]{0.45\textwidth}
        \centering
        \includegraphics[width=\textwidth]{images/Logo - colour PNG.png}
        \caption{Caption A.}
        \label{fig:captiona}
    \end{subfigure}
    \hfill
    \begin{subfigure}[h]{0.45\textwidth}
        \centering
        \includegraphics[width=\textwidth]{images/Logo - colour PNG.png}
        \caption{Caption B.}
        \label{fig:captionb}
    \end{subfigure}
    \caption{Bristol Logo.}
\end{figure}

\subsection{Figures side by side}

\begin{figure}[h]
    \centering
    \begin{minipage}[h]{0.45\textwidth}
        \centering
        \includegraphics[width=\textwidth]{images/Logo - colour PNG.png}
        \caption{Caption C.}
        \label{fig:captionc}
    \end{minipage}
    \hfill
    \begin{minipage}[h]{0.45\textwidth}
        \centering
        \includegraphics[width=\textwidth]{images/Logo - colour PNG.png}
        \caption{Caption D.}
        \label{fig:captiond}
    \end{minipage}
\end{figure}


%%%%%%%%%%%
% Tables  
%%%%%%%%%%%
\section{Tables}

Example table (see Table~\ref{tb:table-example}). 
% Generate LaTeX table quickly with this website: https://www.tablesgenerator.com

\begin{table}[h]
\caption{Table caption.}
\label{tb:table-example}
\begin{tabularx}{\textwidth}{XXX}
\toprule
\textbf{Col 1} & \textbf{Col 2} & \textbf{Col 3} \\ \midrule
\textbf{A}     & 1              & 5              \\
\textbf{B}     & 2              & 6              \\
\textbf{C}     & 3              & 7              \\
\textbf{D}     & 4              & 8              \\ \bottomrule
\end{tabularx}%
\end{table}


%%%%%%%%%%%
% Lists  
%%%%%%%%%%%
\section{Lists}

\subsection{Numbered list}

\subsubsection{Enumeration}

\begin{enumerate}
    \item First item.
    \item Second item.
    \item Third item.
\end{enumerate}

\subsubsection{Research questions}

\begin{enumerate}[font={\bfseries},label={RQ\arabic*}.,align=left]
    \item \rqone
    \item \rqtwo
    \item \rqthree
\end{enumerate}

\subsubsection{Hypothesis}

\begin{enumerate}[font={\bfseries},label={H\arabic*}.,align=left]
    \item First hypothesis.
    \item Second hypothesis.
    \item Third hypothesis.
\end{enumerate}

\subsection{Itemised list}

\begin{itemize}
    \item First item.
    \item Second item.
    \item Third item.
\end{itemize}

\section{Equations}
\begin{equation}
\label{eq:CR}
CR = \frac{(\sum_{i=1}^n \lambda_i)^2}{(\sum_{i=1}^n \lambda_i)^2+(\sum_{i=1}^n \delta_i)} 
\end{equation}
\begin{equation}
\label{eq:AVE}
AVE = \frac{\sum_{i=1}^n \lambda_i^2}{n}
\end{equation}
\begin{equation}
\label{eq:CA}
CA = \frac{n}{n-1}(1-\frac{\sum_{i=1}^n V_i}{V_t}) 
\end{equation}

